\subsubsection{Fourier Transforms of Tempered Distributions}\label{subsec:tempered-fourier-transform}

The previous theorem motivates the definition of the Fourier transform of a tempered distribution.

\begin{Definition}
\label{def:tempered-fourier-transform}
Let $f\in\mathcal{S}^{\prime}(\mathbb{R}^m)$. Then the \textbf{Fourier transform} of $f$ is defined by the functional
\begin{equation}\label{eq:tempered-distribution-fourier-definition}
(\mathcal{F}[f],\phi)=(f,\mathcal{F}^{-1}[\phi]),\quad\phi\in\mathcal{S}(\mathbb{R}^m).
\end{equation}

\end{Definition}

It is clear from the definition that $\mathcal{F}$ is a continuous mapping from $\mathcal{S}^{\prime}(\mathbb{R}^m)$ into itself. It is also easy to check that the formulas \eqref{eq:fourier-frequency-derivatives} and \eqref{eq:fourier-spatial-derivatives} still hold in $\mathcal{S}^{\prime}(\mathbb{R}^m)$; the same is true for the inversion formula \eqref{eq:fourier-inversion-formula}, which can be restated as
\begin{equation}\label{eq:fourier-transform-twice-reflection}
\mathcal{F}\mathcal{F}[f(\bm{x})]=f(-\bm{x});
\end{equation}
this form has meaning for generalized functions.

We shall now consider a number of examples.

\begin{Example}
\label{ex:fourier-delta}
We have
\begin{equation}\label{eq:fourier-transform-delta-pairing}
\begin{aligned}
(\mathcal{F}[\delta],\phi)&=(\delta,\mathcal{F}^{-1}[\phi])=\mathcal{F}^{-1}[\phi](0)\\
&=(2\pi)^{-m/2}\int_{\mathbb{R}^m}\phi(\bm{x})\,\mathrm{d}\bm{x},
\end{aligned}
\end{equation}
i.e., the Fourier transform of $\delta$ is the constant $(2\pi)^{-m/2}$.
\end{Example}

\begin{Example}
\label{ex:fourier-constant}
We have
\begin{equation}\label{eq:fourier-transform-constant-pairing}
\begin{aligned}
(\mathcal{F}[1],\phi)&=(1,\mathcal{F}^{-1}[\phi])=\int_{\mathbb{R}^m}\mathcal{F}^{-1}[\phi](\bm{x})\,\mathrm{d}\bm{x}\\
&=(2\pi)^{m/2}\mathcal{F}\mathcal{F}^{-1}[\phi](0)=(2\pi)^{m/2}\phi(0),
\end{aligned}
\end{equation}
i.e., the Fourier transform of 1 is $(2\pi)^{m/2}\delta$. From \eqref{eq:fourier-frequency-derivatives}, \eqref{eq:fourier-spatial-derivatives}, it is now clear that the Fourier transforms of polynomials are linear combinations of derivatives of the delta function and vice versa.
\end{Example}

\begin{Example}
\label{ex:fourier-exponential-periodic}
A calculation similar to that in Example~\ref{ex:fourier-constant} shows that the Fourier transform of $\exp(i\bm{\eta}\cdot\bm{x})$ (viewed as a function of $\bm{x}$) is $(2\pi)^{m/2}\delta(\bm{\xi}-\bm{\eta})$. If $f$ is a periodic distribution given by a Fourier series
\begin{equation}\label{eq:periodic-distribution-fourier-series}
f(x)=\sum_{n=-\infty}^{\infty}c_ne^{inx},
\end{equation}
we find that
\begin{equation}\label{eq:periodic-distribution-fourier-point-masses}
\mathcal{F}[f](\xi)=\sqrt{2\pi}\sum_{n=-\infty}^{\infty}c_n\delta(\xi-n).
\end{equation}

\end{Example}

\begin{Example}
\label{ex:fourier-compact-distribution-entire}
Let $f$ be a distribution with compact support. Then we can define $(f,\phi)$ for any $\phi\in C^\infty(\mathbb{R}^m)$: we set $(f,\phi)=(f,\phi_0)$, where $\phi_0$ is any element of $\mathcal{D}(\mathbb{R}^m)$ which agrees with $\phi$ in a neighborhood of the support of $f$. It follows from the definition of the support that this definition does not depend on the choice of $\phi_0$. In particular, this defines $f$ as an element of $\mathcal{S}^{\prime}(\mathbb{R}^m)$. We claim now that $\mathcal{F}[f]$ is the function
\begin{equation}\label{eq:compact-distribution-fourier-entire-formula}
\mathcal{F}[f](\bm{\xi})=(2\pi)^{-m/2}(\overline{f(\bm{x})},e^{-i\bm{\xi}\cdot\bm{x}}).
\end{equation}
Here $(\overline f,\phi)$ is defined as the complex conjugate of $(f,\overline\phi)$. To verify the claim, we must show that, for any $\phi\in\mathcal{S}(\mathbb{R}^m)$, we have
\begin{equation}\label{eq:compact-distribution-fourier-integral-interchange}
\begin{aligned}
&(2\pi)^{-m/2}\int_{\mathbb{R}^m}(f(\bm{x}),e^{i\bm{\xi}\cdot\bm{x}})\phi(\bm{\xi})\,\mathrm{d}\bm{\xi}\\
&\quad=(f,\mathcal{F}^{-1}[\phi])\\
&\quad=(2\pi)^{-m/2}\left(f,\int_{\mathbb{R}^m}e^{i\bm{\xi}\cdot\bm{x}}\phi(\bm{\xi})\,\mathrm{d}\bm{\xi}\right).
\end{aligned}
\end{equation}
That is, we have to justify taking $f$ under the integral, which is accomplished in the usual way by approximating the integral by finite sums. We note that \eqref{eq:compact-distribution-fourier-entire-formula} defines an entire function of $\bm{\xi}\in\mathbb{C}^m$. The fact that a distribution of compact support has finite order (Lemma~\ref{lem:distribution-local-finite-order}) implies that for real arguments this function has polynomial growth.
\end{Example}
